\documentclass[10pt,twocolumn]{article}
\usepackage[a4paper,margin=0.72in,columnsep=0.28in]{geometry}
\usepackage[T1]{fontenc}
\usepackage[utf8]{inputenc}
\usepackage{amsmath,amsthm}
\usepackage{newtxtext,newtxmath}
\usepackage{microtype}
\usepackage{booktabs}
\usepackage{xcolor}
\usepackage{enumitem}
\usepackage[numbers,sort&compress]{natbib}
\usepackage[hidelinks]{hyperref}
\usepackage[capitalise,noabbrev]{cleveref}
\usepackage{caption}
\captionsetup{font=small,labelfont=bf,skip=4pt}
\setlist{nosep,leftmargin=*}

\newtheorem{proposition}{Proposition}
\newtheorem{lemma}{Lemma}
\newtheorem{theorem}{Theorem}
\theoremstyle{remark}
\newtheorem*{remark}{Remark}

\newcommand{\F}{\mathbb{F}_p}
\newcommand{\au}[1]{[#1]}
\newcommand{\Fzk}{\mathcal{F}_{\mathsf{ZK}}}
\newcommand{\Flk}{\mathcal{F}^{\mathsf{Lookup}}_{\mathsf{ZK}}}
\newcommand{\RtoF}{\mathsf{R2F}}
\newcommand{\FtoR}{\mathsf{F2R}}
\newcommand{\DD}{\mathsf{DigitDec}}
\newcommand{\Rng}{\mathsf{Range}}
\newcommand{\Lk}{\mathsf{Lookup}}
\newcommand{\CZ}{\mathsf{CheckZero}}
\newcommand{\Trunc}{\mathsf{Trunc}}
\newcommand{\PTrunc}{\mathsf{PosTrunc}}
\newcommand{\Msnzb}{\mathsf{Msnzb}}
\newcommand{\ulp}{\Delta}

% protocol boxes
\newcommand{\protobox}[2]{\noindent\fbox{\begin{minipage}{\dimexpr\columnwidth-2\fboxsep-2\fboxrule\relax}\small
  \centering\textbf{#1}\par\smallskip\raggedright #2\end{minipage}}}

\newcommand{\SpeedMin}{4.0}
\newcommand{\SpeedMax}{7.1}
\newcommand{\CommMin}{4.1}
\newcommand{\CommMax}{7.1}
\newcommand{\PolySpeedMin}{1.6}
\newcommand{\PolySpeedMax}{3.5}
\newcommand{\OverMin}{7}
\newcommand{\OverMax}{16}
\newcommand{\PlainMin}{6}
\newcommand{\PlainMax}{29}
\newcommand{\SinGapLow}{6.1}
\newcommand{\SinGapHigh}{7.5}
\newcommand{\CrossN}{35{,}000}
\newcommand{\VOneLogExOut}{50.40}
\newcommand{\VOneLogExTrue}{26.34}
\newcommand{\VOneSinErr}{46.5}
\newcommand{\Setup}{All runs use one AMD Ryzen 7 9800X3D 8-Core Processor machine with Linux 7.1, a Debian~12 container, GCC~12 and \texttt{-O3}.}
\newcommand{\SoundTotal}{878}
\newcommand{\HaoCommDev}{0.1\%}
\newcommand{\HaoGeluDev}{28\%}
\newcommand{\HaoTimeFactor}{4--5$\times$}
\newcommand{\ExtLines}{939}


\title{\Large\bfseries Sound Zero-Knowledge Proofs for Logarithmic and Periodic\\ Activations in a Lookup-Based Framework}
\author{Zhi Chen Xiang}
\date{}

\begin{document}
\maketitle

\begin{abstract}
Lookup tables are the key tool behind the most efficient VOLE-based zero-knowledge (ZK) proofs for
non-linear functions in machine learning~\cite{hao2024scalable}. A natural way to extend such a
framework is to reuse its pattern for the exponential: split the input into digits, look up each
digit, and recombine the results. An earlier master's thesis~\cite{chen2025thesis} sketched a
logarithm and a sine in this way, which would enable the softplus, log-softmax and snake
activations, but did not implement them. Implementing them now shows that the template needs care.
The digit-wise recombination does not compute the logarithm and needs undefined table entries on
typical inputs, and the sine's range reduction must be bound to the input for the proof to be
sound. We isolate these two pitfalls and give revised protocols that use only the framework's
building blocks: a logarithm normalised with $\Msnzb$, and a protocol for periodic functions that performs a
\emph{verified} reduction modulo one. We prove error bounds of $1$--$1.5$ units in the last place
($\ulp=2^{-12}$). We implement both protocols in Hao et al.'s released C++ code base (emp-zk). The implementation matches an
executable specification bit for bit and rejects every tampered witness. For $10^5$ instances, the
lookup-based protocols are \SpeedMin--\SpeedMax$\times$ faster and use
\CommMin--\CommMax$\times$ less communication than the same functions proven without lookups in
the same backend.
\end{abstract}

% =====================================================================================
\section{Introduction}

ZK proofs let a model owner prove that an inference was computed correctly without revealing
the weights. In practice, the non-linear layers are the bottleneck. Hao et al.~\cite{hao2024scalable}
remove most of this cost with \emph{table lookups} over a VOLE/IT-MAC backend. Inputs are split into
12-bit digits with a verified digit decomposition, and each digit indexes a small public table.
Their framework covers ReLU, sigmoid, GELU, softmax and normalisation, and reports
50--179$\times$ speed-ups over Mystique~\cite{weng2021mystique}.

The author's master's thesis~\cite{chen2025thesis} (``v1'' below) proposed extending the framework
with a logarithm $\Pi_{\log}$ and a sine $\Pi_{\sin}$. These are needed for softplus, log-softmax
and the periodic snake activation~\cite{ziyin2020snake}, which is used in neural audio
codecs~\cite{kumar2023dac}. Both protocols follow the template of the exponential $\Pi_{\exp}$: a
digit decomposition, one lookup per digit, and a cheap recombination. The thesis presented them as
paper designs, without an implementation to check them against.

\paragraph{This paper.} We implement both protocols and check them in code. This shows that the
template does not carry over directly, and we revise the designs accordingly.
\begin{enumerate}
\item \textbf{Two pitfalls} (\cref{sec:pitfalls}). (i) Digit-wise lookups only compose through a
homomorphism: $e^{-\sum c_i}=\prod e^{-c_i}$, but $\log\sum c_i\neq\sum\log c_i$. The v1 logarithm
outputs $\log\prod_i c_i$, and it needs the undefined entry $\log 0$ for every input below
$2^{36}$. (ii) A committed intermediate that no constraint ties to the input makes a protocol
unsound. v1's sine does not yet tie the reduced angle $x_{\bmod}$ to the input, so a malicious
prover can have any table output accepted.
\item \textbf{Sound constructions} (\cref{sec:constructions}). $\Pi_{\mathsf{Per}}$ evaluates any
1-periodic function in \emph{turns}. The reduction mod~1 is exactly a digit decomposition, so it is
verified at no extra cost, and one $2^{12}$-row table with first-order interpolation reaches
$\approx\!1\ulp$. $\Pi_{\log}$ normalises multiplicatively with $\Msnzb$, as Hao's division does,
and reaches $1.5\ulp$ over $(0,2^{60})$. Both come with error bounds and a security argument in the
same hybrid model as~\cite{hao2024scalable}. We specialise them to snake, softplus and log-softmax.
\item \textbf{Implementation and evaluation} (\cref{sec:impl,sec:eval}). We add the protocols to
Hao et al.'s released C++ code (ZKMath on emp-zk~\cite{emptoolkit}). A pure-Python executable
specification comes with an adversarial prover, and the C++ output matches it bit for bit. On the
real backend, every tampered witness is rejected. The benchmarks follow Hao et al.'s methodology:
$10^3$--$10^5$ instances and 200\,Mbps--1\,Gbps links. We compare against the same functions
proven \emph{without} lookups in the same backend, and against plaintext.
\end{enumerate}

\paragraph{What is (not) new.} Evaluating $\log$ or $\sin$ through lookups is not new in itself.
zkLLM's \textsf{tlookup} handles arbitrary element-wise functions and uses a logarithm inside
zkAttn~\cite{sun2024zkllm}. Lookup-based non-linearities are standard in SNARK-based
zkML~\cite{kang2022scaling}. Normalisation plus table lookup is the usual approach in secure
computation~\cite{rathee2021sirnn,aly2019benchmarking}. SwiGLU is directly expressible in
Hao et al.'s framework and in zkLLM. Our contribution is to identify the pitfalls, to give sound
and precise instantiations \emph{inside} the digit-decomposition framework together with their
proofs, and to provide an open, reproducible implementation and benchmark.

% =====================================================================================
\section{Background}\label{sec:bg}

\paragraph{Fixed point.} We use Hao et al.'s parameters: $p=2^{61}-1$ and $s=12$ fractional bits. A
real $\hat x$ is encoded as $x=\RtoF(\hat x)=\lfloor\hat x\cdot2^s\rceil\bmod p$, and decoded as
$\hat x=\FtoR(x)=x/2^s$, with $x$ read as a signed integer in $(-p/2,p/2)$. Errors are stated in
units of $\ulp=2^{-s}\approx2.44\cdot10^{-4}$. Public table entries are rounded to nearest
(error $\le\ulp/2$).

\paragraph{VOLE-based ZK.} The prover $P$ commits to values as IT-MACs $\au{x}$. Linear operations
with public constants are free. A multiplication $\au{x}\au{y}$ costs one committed value and is
checked in batches as in QuickSilver~\cite{yang2021quicksilver}. $\CZ$ opens a zero-test. We work in
the $(\Fzk,\Flk)$-hybrid model of~\cite{hao2024scalable}. $\Lk(L,\au{x},\au{\vec y})$ checks that
$(x,\vec y)$ is a row of a public table $L$, and $\Rng(d,\au{x})$ checks that $x\in[0,2^d)$. Both are
realised with a ZK read-only memory~\cite{franzese2021ram}. $N$ lookups into a $T$-row table cost
$2N+T$ multiplications, i.e.\ about two per lookup, plus one committed version per read.

\paragraph{Building blocks~\cite{hao2024scalable}.}
$\DD(\au{x};d_0,\dots,d_{k-1})$ commits digits $x_i$, range-checks each of them (digits of up to
12 bits use one lookup), and runs $\CZ(\sum_i2^{\sigma_i}x_i-x)$ with $\sigma_i=\sum_{j<i}d_j$.
Because $\sum d_i\le60<\log p$, an accepted decomposition is unique and proves $x<2^{\sum d_i}$.
$\PTrunc(\au{x},t)$ is the high digit of $\DD(\au{x};t,n-t)$, optionally after adding $2^{t-1}$
for rounding. We compute a signed truncation for $|x|<2^L$ as $\PTrunc(x+2^L,t)-2^{L-t}$, which
costs one $\DD$. $\Msnzb$ proves $2^k\le x<2^{k+1}$: $P$ commits $k$ and the row $(k,2^k,\dots)$ of
a public table, and $V$ range-checks $x-2^k$ and $2^{k+1}-1-x$ on $n$ bits. For $x<2^n$ this is
equivalent to Hao et al.'s comparisons. The exponential $\Pi_{\exp}$
computes $e^{-\hat x}$ from $\DD(\au{x};12,12,\dots)$ with one table per digit,
$L_i=\{v\mapsto e^{-2^{\sigma_i}v/2^s}\}$, and multiplies the looked-up factors with truncation.

% =====================================================================================
\section{Two pitfalls}\label{sec:pitfalls}

\subsection{Digit-wise lookups need a homomorphism}
$\Pi_{\exp}$ is correct because the digits of $x$ split its value additively,
$\hat x=\sum_i c_i$ with $c_i=2^{\sigma_i}\hat x_i$, and the exponential maps sums to products.
v1's $\Pi_{\log}$~\cite[Fig.~4.1]{chen2025thesis} copies the structure. It looks up
$y_i=\RtoF(\log c_i)$ in per-digit tables and outputs $\sum_i y_i$, justified by
``the logarithm of a product is the sum of the logarithms''. The digits, however, combine as a sum.

\begin{proposition}\label{prop:log}
With exact tables, v1's $\Pi_{\log}$ outputs $\log\prod_ic_i$ instead of $\log\hat x=\log\sum_ic_i$,
an error of $\log(\prod_ic_i/\sum_ic_i)$ whose size and sign depend on the input. The table entry
for a zero digit would be $\log 0$ and does not exist, so every input with a zero digit makes the
honest prover abort. Hence, for $k\ge2$ digits, the protocol returns $\ln\hat x$ only by coincidence.
With Hao et al.'s 12-bit digits ($k=5$ for $x<2^{60}$), every $x<2^{48}$, i.e.\ every
$\hat x<2^{36}$, has a zero top digit and aborts.
\end{proposition}
For example, $x=2^{50}+2^{40}+3\cdot2^{24}+5\cdot2^{12}+7$ has five non-zero digits. v1 outputs
\VOneLogExOut\ while $\ln\hat x=\VOneLogExTrue$. Every probability, every softmax denominator
in $[1,n]$ and every softplus argument in $[1,2]$ lies far below $2^{36}$, so the honest prover
cannot complete the protocol on the inputs it was designed for (last column of \cref{tab:prec}).
The thesis text also writes ``$k=12$'', which is the digit size $d$ rather than the number of
digits $k$.

\subsection{Every witness must be pinned to the input}
v1's $\Pi_{\sin}$~\cite[Fig.~4.2]{chen2025thesis} works as follows. $P$ computes
$x_{\bmod}=\hat x-2\pi\lfloor\hat x/2\pi\rfloor$ locally and commits $\au{x_{\bmod}}$. It then commits
$y=\RtoF(T_5(x_{\bmod}))$ with the Taylor polynomial $T_5(t)=t-t^3/6+t^5/120$, and the parties
run $\Lk(L,\au{x_{\bmod}},\au{y})$.

\begin{proposition}\label{prop:sin}
v1's $\Pi_{\sin}$ is not sound. For every input $\au{x}$ and every row $(t,y)$ of $L$, the prover
that commits $x_{\bmod}=t$ is accepted with output $y$.
\end{proposition}
\begin{proof}
$\au{x}$ does not appear in any check. The only check involving $x_{\bmod}$ is the lookup, which
passes for every row.
\end{proof}
Our test suite runs this attack against a literal implementation of v1 and obtains an accepted
proof of $\sin(0)=1.0044$. Running the honest protocol also shows larger error margins than the
thesis estimated. On the protocol's domain $[0,2\pi)$, $T_5$ is off by up to \VOneSinErr\ at
$t\to2\pi^-$, and the degree-7 polynomial that the thesis suggests for a full cycle still errs by
30.2 there. Negative inputs wrap to $x_{\bmod}\approx2\pi$, so $\hat x\in[-1,0)$ already gives errors
of 46. Moreover, a lookup-based protocol does not need a Taylor polynomial at all: the table is
public, so it can store $\sin$ itself. The real difficulty is the \emph{verified} range reduction.
It cannot be exact modulo $2\pi$, because $2\pi$ is irrational and a fixed-point constant drifts by
$q\cdot|C2^{-s}-2\pi|$ after $q$ periods.

\paragraph{Lessons.} (L1) Every committed value must be a function of the inputs, fixed by $\CZ$,
range and lookup relations. The soundness proofs in the hybrid model, including all of Hao et al.'s,
rely on this \emph{witness uniqueness}. (L2) A digit-wise lookup recombines
correctly only if $f(\sum_ic_i)=\bigodot_if(c_i)$ for a cheap operation $\bigodot$. This holds for
$\exp$ with $\bigodot=\times$. It also holds for periodic functions whose period, in fixed-point
units, is a power of two: the digits above the period contribute whole periods and drop out. For
$\log$ one needs a multiplicative decomposition, i.e.\ normalisation.

% =====================================================================================
\section{Sound constructions}\label{sec:constructions}

\begin{figure}[t]
\protobox{Protocol $\Pi_{\mathsf{Per}}[g]$ (1-periodic $g$, e.g.\ $\sin2\pi\tau$, $\sin^22\pi\tau$)}{%
\textbf{Parameters:} $b=12$; table $L_g=\{(i,\RtoF(g(i2^{-b})),\RtoF(g'(i2^{-b})))\}_{i<2^b}$.\\
\textbf{Input:} $\au{u}$ with $|u|<2^{58}$, read as $\tau=u/2^F$, $b\le F\le48$.
\begin{enumerate}
\item $\au{u'}=\au{u}+2^{58}$. Since $2^F\mid2^{58}$, $\mathrm{frac}(u'/2^F)=\mathrm{frac}(\tau)$.
\item $(\au{\delta},\au{i},\au{h})\leftarrow\DD(\au{u'};F-b,\,b,\,59-F)$. This is the verified
  reduction mod 1: $\mathrm{frac}(\tau)=i2^{-b}+\delta2^{-F}$; the top digit proves $|u|<2^{58}$.
\item $P$ commits $\au{g_0},\au{g_1}$; $\Lk(L_g,\au{i},\au{g_0},\au{g_1})$.
\item $\au{y}=\au{g_0}+\Trunc(\au{g_1}\cdot\au{\delta},\,F)$ (first-order interpolation).
\end{enumerate}}
\caption{Evaluation of periodic functions in turns.}\label{fig:per}
\end{figure}

\begin{figure}[t]
\protobox{Protocol $\Pi_{\log}$}{%
\textbf{Parameters:} $n\le60$, $m=12$; tables $K=\{(k,2^k,2^{n-1-k},\RtoF((k-s)\ln2))\}_{k<n}$ and
$M=\{(z_1,\RtoF(\ln(z_12^{-m})),\RtoF(2^m/z_1))\}_{2^m\le z_1<2^{m+1}}$.\\
\textbf{Input:} $\au{x}$ with $x\in(0,2^n)$; output $\approx\ln\hat x$.
\begin{enumerate}
\item ($\Msnzb$) $P$ commits $\au{k},\au{e},\au{q},\au{c}$; $\Lk(K,\au{k},\au{e},\au{q},\au{c})$;
  $\Rng(n,\au{x}-\au{e})$; $\Rng(n,2\au{e}-1-\au{x})$.
\item $\au{z}=\au{x}\cdot\au{q}$, so $z\in[2^{n-1},2^n)$ (multiplicative normalisation).
\item $(\au{z_0},\au{z_1})\leftarrow\DD(\au{z};n-1-m,\,m+1)$; $P$ commits $\au{a},\au{b}$;
  $\Lk(M,\au{z_1},\au{a},\au{b})$.
\item $\au{y}=\au{c}+\au{a}+\PTrunc_{\mathrm{round}}(\au{b}\cdot\au{z_0},\,n-1)$.
\end{enumerate}}
\caption{Logarithm via normalisation (cf.\ Hao et al.'s $\Pi_{\mathsf{Div}}$).}\label{fig:log}
\end{figure}

\begin{table}[t]\centering\footnotesize
\caption{Accuracy in units of $\Delta=2^{-12}$ over $10^5$ random inputs (max\,/\,mean absolute error); cost per instance of \textsf{Ours} (function lookups\,+\,range checks\,/\,multiplications; log-softmax per element); and the v1 protocols~\cite{chen2025thesis} with an honest prover (max error or share of aborted runs).}
\label{tab:prec}
\setlength{\tabcolsep}{2.5pt}
\begin{tabular}{lcccl}\toprule
Function & \textsf{Ours} & \textsf{Naive} & Cost & v1 \\\midrule
$\sin x$ & 1.00\,/\,0.33 & 0.53\,/\,0.25 & 1+10\,/\,1 & err 46\\
$\ln x$, $n{=}60$ & 1.49\,/\,0.40 & 1.03\,/\,0.33 & 2+22\,/\,2 & err 49; abort 79\%\\
$\ln x$, $n{=}20$ & 1.46\,/\,0.36 & 1.01\,/\,0.30 & 2+10\,/\,2 & abort 100\%\\
snake$_1$ & 1.00\,/\,0.33 & 0.54\,/\,0.25 & 1+14\,/\,1 & err $\le$\,2166\\
softplus & 1.19\,/\,0.18 & 0.98\,/\,0.15 & 4+18\,/\,4 & abort 100\%\\
log-softmax$_{10}$ & 3.82\,/\,0.54 & 2.99\,/\,0.46 & 2.2+8\,/\,2.1 & abort 100\%\\
\bottomrule\end{tabular}
\end{table}


\subsection{Periodic functions in turns}
We parametrise periodic functions by \emph{turns}. We write $\sin(a\hat x)=\sin(2\pi\tau)$ with
$\tau=a\hat x/2\pi$, and fold the constant $a/2\pi$ into a public multiplier. For snake, $a$ is a
trained weight, so the constant can be folded in at training time. Concretely,
$u=x\cdot A$ with $A=\mathrm{round}(a2^{s_a}/2\pi)$ is a local operation and gives
$\tau=u/2^{s+s_a}$. The reduction mod 1 is then the low $F$ bits of $u$, which is exactly what $\DD$
extracts. \Cref{fig:per} gives the protocol.

\begin{lemma}\label{lem:per}
For $g\in C^2$ with period 1, $\Pi_{\mathsf{Per}}$ outputs $y$ with
$|\FtoR(y)-g(\tau)|\le\ulp+2^{-s-1-b}+\tfrac12\max|g''|\,2^{-2b}$.
For $g=\sin2\pi\tau$ and $b=12$ this is $1.005\ulp$. The conversion from radians adds at most
$2\pi|\hat x|\,|A2^{-s_a}-a/2\pi|\le2\pi|\hat x|2^{-s_a-1}$, which is $1.9\cdot10^{-7}$ for
$|\hat x|<4096$ and $s_a=36$.
\end{lemma}
\begin{proof}
Write $\tau_i=i2^{-b}$ and $\hat\delta=\delta2^{-F}\in[0,2^{-b})$. The table entries carry errors
$|e_0|,|e_1|\le\ulp/2$, and the rounded truncation adds $\ulp/2$. The term $e_1\hat\delta$ is
below $2^{-s-1-b}$. The Taylor remainder of $g(\tau_i+\hat\delta)$ around $\tau_i$ is at most
$\max|g''|\hat\delta^2/2$.
\end{proof}
\noindent\emph{Cost:} one lookup into a $2^{12}$-row table, one multiplication, and the range
lookups of two $\DD$s (59 and $F+5$ bits; 10 range lookups for $F=48$, see \cref{tab:prec}).
Snake, $\mathrm{snake}_a(x)=x+\sin^2(ax)/a$, tabulates $g=\sin^2(2\pi\tau)$ directly, so no
squaring is needed, and then multiplies by $1/a$ (a 24-bit constant) followed by one truncation.
We treat $a$ as public: it is fixed per channel after training. If $a$ must stay private, $a/2\pi$
and $1/a$ are committed as model parameters, which costs one extra multiplication per activation.

\subsection{Logarithm via normalisation}
\Cref{fig:log} normalises exactly like the first steps of Hao et al.'s division. $\Msnzb$ gives $k$
with $2^k\le x<2^{k+1}$, and multiplying by the looked-up $2^{n-1-k}$ yields
$\hat z=z/2^{n-1}\in[1,2)$ with $\ln\hat x=(k-s)\ln2+\ln\hat z$. The top $m+1$ bits $z_1$ of $z$
index the mantissa table, and the remainder $r=z_0/2^{n-1}<2^{-m}$ is handled by
$\ln(\hat z_1+r)=\ln\hat z_1+r/\hat z_1-\theta$, where $0\le\theta\le2^{-2m-1}$.

\begin{lemma}\label{lem:log}
For every $x\in(0,2^n)$, $|\FtoR(y)-\ln\hat x|\le\tfrac32\ulp+2^{-2m-1}+2^{-m-s-1}$, which is
$3.66\cdot10^{-4}$ ($1.5\ulp$) for $m=12$.
\end{lemma}
\begin{proof}
The entries $c$ and $a$ and the rounding of $\PTrunc$ each contribute at most $\ulp/2$. The error
$e_b$ of $b$ enters as $r\,e_b\le2^{-m-s-1}$, and the linearisation contributes $\theta$.
\end{proof}
\noindent\emph{Cost:} 2 lookups (into $n$- and $2^{12}$-row tables), 2 multiplications, and
$2\lceil n/12\rceil$ range checks for $\Msnzb$ plus those of two $\DD$s: 22 range checks for
$n=60$ and 10 for $n=20$, two of them single-bit (\cref{tab:prec}). $M$ contains no row with $z_1<2^m$, so the lookup
itself enforces the normalisation. The protocol aborts on $x\le0$ and $x\ge2^n$.

\subsection{Activations}
\textbf{Softplus.} We compute $\mathrm{softplus}(x)=\mathrm{ReLU}(x)+\ln(1+e^{-|x|})$. A
sign/absolute-value step costs one bit, one multiplication and one range check. $\Pi_{\exp}$ then
evaluates $e^{-|x|}$, and $\Pi_{\log}$ runs with $n=s+2$ because the argument lies in $[1,2]$.
\textbf{Log-softmax.} We compute $y_j=(x_j-m)-\ln\sum_ke^{x_k-m}$, where $m$ comes from Hao et al.'s
max (range checks plus $\prod_j(m-x_j)=0$). The sum lies in $[1,n]$, so $\Pi_{\log}$ runs on
$s+\lceil\log_2n\rceil+1$ bits, and no division is needed, unlike softmax.
\textbf{Snake:} see above. Swish, GLU and SwiGLU need only Hao et al.'s sigmoid and a
multiplication, and are not claimed here.

\subsection{Security}
We define the ideal functionality $\mathcal{F}^{+}_{\mathsf{Math}}$ as Hao et al.'s
$\mathcal{F}_{\mathsf{Math}}$ extended with commands $\mathsf{Log}$, $\mathsf{Per}_g$, $\mathsf{Snake}$,
$\mathsf{Softplus}$ and $\mathsf{LogSoftmax}$. Each command outputs the deterministic fixed-point
function computed by the honest protocol, just as~\cite{hao2024scalable} defines
$\mathsf{PtExp}$, and aborts outside the domain. Lemmas~\ref{lem:per} and~\ref{lem:log} then bound
the distance to the real functions.

\begin{theorem}
$\Pi_{\mathsf{Per}}$, $\Pi_{\log}$ and the activation protocols UC-realise~\cite{canetti2001uc}
$\mathcal{F}^{+}_{\mathsf{Math}}$ in the $(\Fzk,\Flk)$-hybrid model against a static, malicious
adversary.
\end{theorem}
\begin{proof}[Proof sketch]
\emph{Corrupted verifier:} $V$ receives only accept/abort bits from the hybrid functionalities and
never an opened value, so the simulator runs $V$ on dummy inputs. \emph{Corrupted prover:} the
simulator extracts every committed witness from the $\mathsf{Input}$ calls and replays the checks.
If all checks pass, the witnesses are the honest ones (L1). The digits of each $\DD$ are unique
because $\sum d_i\le60<\log p$. $\Msnzb$'s range checks fix $k$. Every table is a function of its
key column. Products and linear combinations are determined. Hence the output equals the ideal
output, and the simulator sends the extracted input to $\mathcal{F}^{+}_{\mathsf{Math}}$. The
structure follows the proofs in~\cite[App.~B]{hao2024scalable}. We cross-validate the argument
empirically in \cref{sec:impl}.
\end{proof}

\begin{table*}[t]\centering\small
\caption{Runtime (s) and communication (MB) for $10^5$ instances (log-softmax: $10^5$ rows of 10), in the format of Hao et al.'s Tables 4--5. Speed-ups of \textsf{Ours} over the baseline in parentheses. \textsf{Plain}: floating-point evaluation without proof.}
\label{tab:main}
\begin{tabular}{llrrrrr}\toprule
Function & Protocol & 200\,Mbps & 500\,Mbps & 1\,Gbps & Comm.\ (MB) & Elem./inst.\\\midrule
$\sin x$ & \textsf{Ours} & 1.610 & 0.663 & 0.350 & 37.2 & 49\\
 & \textsf{Poly} & 4.758 (3.0$\times$) & 2.144 (3.2$\times$) & 1.325 (3.8$\times$) & 108.2 (2.9$\times$) & 142\\
 & \textsf{Naive} & 9.777 (6.1$\times$) & 4.313 (6.5$\times$) & 2.627 (7.5$\times$) & 222.5 (6.0$\times$) & 292\\
 & \textsf{Plain} & \multicolumn{3}{c}{1.64 ms} & -- & --\\
\midrule
$\ln x$, $n{=}60$ & \textsf{Ours} & 3.522 & 1.629 & 1.053 & 79.3 & 104\\
 & \textsf{Poly} & 6.190 (1.8$\times$) & 2.647 (1.6$\times$) & 1.515 (1.4$\times$) & 142.3 (1.8$\times$) & 187\\
 & \textsf{Naive} & 14.48 (4.1$\times$) & 6.461 (4.0$\times$) & 3.939 (3.7$\times$) & 326.8 (4.1$\times$) & 428\\
 & \textsf{Plain} & \multicolumn{3}{c}{0.84 ms} & -- & --\\
\midrule
$\ln x$, $n{=}20$ & \textsf{Ours} & 1.671 & 0.760 & 0.467 & 38.5 & 50\\
 & \textsf{Poly} & 4.834 (2.9$\times$) & 2.092 (2.8$\times$) & 1.219 (2.6$\times$) & 110.4 (2.9$\times$) & 145\\
 & \textsf{Naive} & 10.13 (6.1$\times$) & 4.500 (5.9$\times$) & 2.679 (5.7$\times$) & 230.2 (6.0$\times$) & 302\\
 & \textsf{Plain} & \multicolumn{3}{c}{1.25 ms} & -- & --\\
\midrule
snake$_1$ & \textsf{Ours} & 2.199 & 0.904 & 0.471 & 51.1 & 67\\
 & \textsf{Poly} & 5.984 (2.7$\times$) & 2.654 (2.9$\times$) & 1.577 (3.4$\times$) & 136.8 (2.7$\times$) & 179\\
 & \textsf{Naive} & 12.70 (5.8$\times$) & 5.631 (6.2$\times$) & 3.364 (7.2$\times$) & 288.4 (5.6$\times$) & 378\\
 & \textsf{Plain} & \multicolumn{3}{c}{2.85 ms} & -- & --\\
\midrule
softplus & \textsf{Ours} & 2.813 & 1.241 & 0.731 & 64.9 & 85\\
 & \textsf{Poly} & 9.598 (3.4$\times$) & 4.258 (3.4$\times$) & 2.547 (3.5$\times$) & 219.2 (3.4$\times$) & 287\\
 & \textsf{Naive} & 19.64 (7.0$\times$) & 8.670 (7.0$\times$) & 5.201 (7.1$\times$) & 446.1 (6.9$\times$) & 585\\
 & \textsf{Plain} & \multicolumn{3}{c}{2.78 ms} & -- & --\\
\midrule
log-softmax$_{10}$ & \textsf{Ours} & 15.51 & 6.910 & 4.159 & 351.2 & 46\\
 & \textsf{Poly} & 53.87 (3.5$\times$) & 23.88 (3.5$\times$) & 14.30 (3.4$\times$) & 1224 (3.5$\times$) & 160\\
 & \textsf{Naive} & 111.0 (7.2$\times$) & 49.19 (7.1$\times$) & 29.36 (7.1$\times$) & 2510 (7.1$\times$) & 329\\
 & \textsf{Plain} & \multicolumn{3}{c}{6.12 ms} & -- & --\\
\bottomrule
\end{tabular}
\end{table*}


% =====================================================================================
\section{Implementation}\label{sec:impl}

\paragraph{C++ on ZKMath.} We extend Hao et al.'s released code (ZKMath, commit
\texttt{b4f03cd}) on emp-zk. The emp-tool/-ot/-zk commits are pinned to December 2024, when ZKMath
was published, and everything builds in a container. Apart from adding our files, the only change to
the original code is the build type: ZKMath's CMake file forces an unoptimised debug build, and we
build everything with \texttt{-O3}. Our extension (\ExtLines\ lines of C++) adds a multi-column
zk-ROM table. Unlike
ZKMath's tables, whose permutation-check challenges come from a fixed PRG key (acceptable for
benchmarking, not for deployment), our tables obtain the challenges from the verifier, at the cost of
one extra round trip per $10^5$ reads. The extension also adds a batched QuickSilver degree-2 check
for bits, which commits no product, and the protocols of \cref{sec:constructions} together with the
baselines of \cref{sec:eval}.

\paragraph{Executable specification.} A pure-Python simulation of the protocols in the
$(\Fzk,\Flk)$-hybrid model serves as the specification. It contains an adaptive malicious prover
that shifts single witnesses, mounts carry attacks on adjacent digits, and tampers with random
pairs of witnesses. Each accepted run must return the honest output. A mutation test confirms that
this fuzzer does find accepted wrong outputs once any single check is disabled.

\paragraph{Cross-validation.} (i) For every function, the C++ protocol run on the real VOLE
backend outputs \emph{exactly} the field elements of the Python specification (\CrossN\ outputs, 0
mismatches). (ii) Malicious prover on the real backend: for each witness $k$ of an execution, the
prover adds 1 to witness $k$. \Cref{tab:sound} shows that all \SoundTotal\ such runs are
rejected, including those against the baselines.

\begin{table}[t]\centering\small
\caption{Malicious prover on the real emp-zk backend: each witness of one execution is incremented in turn. A sound protocol rejects the run or reproduces the honest output; ``wrong'' counts accepted runs with a different output.}
\label{tab:sound}
\setlength{\tabcolsep}{4pt}
\begin{tabular}{llrrr}\toprule
Function & Protocol & Witnesses & Rejected & Wrong\\\midrule
$\sin x$ & \textsf{Ours} & 14 & 14 & 0\\
$\ln x$, $n{=}60$ & \textsf{Ours} & 31 & 31 & 0\\
$\ln x$, $n{=}20$ & \textsf{Ours} & 18 & 18 & 0\\
$e^{-x}$ & \textsf{Ours} & 8 & 8 & 0\\
snake$_1$ & \textsf{Ours} & 20 & 20 & 0\\
softplus & \textsf{Ours} & 29 & 29 & 0\\
log-softmax$_{10}$ & \textsf{Ours} & 119 & 119 & 0\\
$\sin x$ & \textsf{Poly} & 44 & 44 & 0\\
$\ln x$, $n{=}20$ & \textsf{Poly} & 45 & 45 & 0\\
$\sin x$ & \textsf{Naive} & 279 & 279 & 0\\
$\ln x$, $n{=}20$ & \textsf{Naive} & 271 & 271 & 0\\
\bottomrule\end{tabular}
\end{table}


% =====================================================================================
\section{Evaluation}\label{sec:eval}

\paragraph{Setup.} \Setup\ Following~\cite{hao2024scalable}, both parties run single-threaded on
one machine, and the loopback link is shaped to 200\,Mbps, 500\,Mbps and 1\,Gbps. $10^5$ instances
is the default, and the parameters are $p=2^{61}-1$, $s=12$ and 12-bit digits. Our times cover the
whole proof \emph{including} every verifier check: lookup permutations, multiplication triples,
$\CZ$ and degree-2 checks. Communication is the total over both directions.

\paragraph{Baselines.} \textsf{Ours} uses lookups for range checks and for function evaluation.
\textsf{Poly} keeps Hao et al.'s lookup-based building blocks, including $\Msnzb$'s table indexed by
the exponent, but replaces every \emph{function} table by the smallest-degree polynomial with error
below $\ulp/8$ (odd degree~7 for
$\sin2\pi t$ on $[0,\frac14]$ after folding into the first quadrant, degree~5 for $\ln(1+t)$, and
degree~4 for $2^{-f}$ with $e^{-x}=2^{-k-f}$), evaluated by Horner's rule at scale $2^{-20}$.
\textsf{Naive} additionally replaces every lookup-based range check by a bit decomposition with
QuickSilver booleanity checks, and $\Msnzb$ by a prefix-OR over the bits (one multiplication per
bit). It is the ``no table lookup'' version of the same protocols: the
same verified reductions in the same backend, with every lookup removed. The polynomial baselines
are slightly \emph{more} precise than ours (\cref{tab:prec}), so the comparison does not favour
lookups. \textsf{Plain} is the floating-point function in C++.

\paragraph{Calibration.} We re-ran Hao et al.'s own binaries (\cref{tab:hao}). The communication in
both directions matches their published numbers within \HaoCommDev\ for division, sigmoid and
softmax, and for their building blocks: PosTrunc uses 15.57\,MB and Msnzb 49.59\,MB, against 15.571
and 49.586\,MB in their Table~2. The released exponential takes 16-bit inputs, so it uses less
communication than published. The released GELU uses \HaoGeluDev\ more, and we did not investigate
the difference. Our absolute runtimes are \HaoTimeFactor\ lower than the published ones at the same
bandwidth, because of a newer CPU and \texttt{-O3}. We therefore only compare protocols run in the
same setup.

\begin{table}[t]\centering\small
\caption{Calibration: Hao et al.'s released binaries ($10^5$ instances) re-run in our setup vs.\ their published numbers (500\,Mbps).}
\label{tab:hao}
\setlength{\tabcolsep}{4pt}
\begin{tabular}{lrrrr}\toprule
 & \multicolumn{2}{c}{Comm.\ (MB)} & \multicolumn{2}{c}{Time (s), 500\,Mbps}\\\cmidrule(lr){2-3}\cmidrule(lr){4-5}
Protocol & published & ours & published & ours\\\midrule
exp & 99.0 & 28.7 & 8.70 & 0.49\\
div & 110.7 & 110.7 & 9.84 & 2.09\\
sigmoid & 189.9 & 189.9 & 17.71 & 3.64\\
gelu & 338.2 & 432.2 & 32.70 & 8.34\\
softmax & 816.3 & 816.2 & 78.29 & 15.95\\
\bottomrule\end{tabular}
\end{table}


\paragraph{Results.} \Cref{tab:main} gives the results in the format of Hao et al.'s Tables 4--5.
Lookups make the new math functions and activations \SpeedMin--\SpeedMax$\times$ faster than
\textsf{Naive} and use \CommMin--\CommMax$\times$ less communication. Most of the gain already comes
from lookup-based range checks (\textsf{Poly} vs.\ \textsf{Naive}). A 12-bit range check by lookup
costs four committed values (the digit, its version and two permutation-check products), while a
bit decomposition costs twelve. The function tables add a further
\PolySpeedMin--\PolySpeedMax$\times$: they replace Horner steps whose truncations each need a
40-bit decomposition. At these bandwidths all protocols are communication-bound, so the speed-up
closely follows the communication ratio. Only at 1\,Gbps does computation start to matter, and the
gap grows slightly (for $\sin$, from \SinGapLow$\times$ at 200\,Mbps to \SinGapHigh$\times$ at 1\,Gbps). With
lookups, a proof costs \OverMin--\OverMax\,$\mu$s per instance at 500\,Mbps, compared with
\PlainMin--\PlainMax\,ns for the plaintext function. The costs scale linearly with the number of instances, and
the lookup advantage is already present at $10^3$ instances, where the fixed cost of checking the
$2^{12}$-row tables is still visible (\cref{tab:scaling}).

\paragraph{Why not 50--179$\times$?} Hao et al.\ compare against Mystique, which converts
arithmetic values to Boolean ones and evaluates IEEE-754 circuits. We instead compare against the
strongest simple arithmetic alternative in the \emph{same} VOLE backend. The resulting factors of
\SpeedMin--\SpeedMax$\times$ isolate what the lookup technique itself contributes to the new
functions. They should not be read as a comparison with other systems.

\begin{table}[t]\centering\small
\caption{Scaling on the unshaped loopback: runtime (s), median of 3 runs.}
\label{tab:scaling}
\setlength{\tabcolsep}{4pt}
\begin{tabular}{llrrr}\toprule
Function & Protocol & $10^{3}$ & $10^{4}$ & $10^{5}$\\\midrule
$\sin x$ & \textsf{Ours} & 0.006 & 0.032 & 0.210\\
 & \textsf{Naive} & 0.023 & 0.189 & 2.181\\
$\ln x$, $n{=}60$ & \textsf{Ours} & 0.010 & 0.058 & 0.687\\
 & \textsf{Naive} & 0.031 & 0.270 & 3.398\\
$\ln x$, $n{=}20$ & \textsf{Ours} & 0.007 & 0.038 & 0.319\\
 & \textsf{Naive} & 0.023 & 0.197 & 2.278\\
snake$_1$ & \textsf{Ours} & 0.009 & 0.040 & 0.283\\
 & \textsf{Naive} & 0.028 & 0.234 & 2.846\\
softplus & \textsf{Ours} & 0.009 & 0.059 & 0.472\\
 & \textsf{Naive} & 0.039 & 0.353 & 4.375\\
log-softmax$_{10}$ & \textsf{Ours} & 0.032 & 0.212 & 2.803\\
 & \textsf{Naive} & 0.205 & 2.596 & 25.65\\
\bottomrule\end{tabular}
\end{table}


\paragraph{Limitations.} Everything runs in the fixed-point model of~\cite{hao2024scalable}
($s=12$), and our errors are quantisation-limited at $1$--$1.5\ulp$. Accuracy-critical uses need a
larger $s$ and larger tables. The proofs are for the protocols in the hybrid model. The C++ code is a
research prototype built on an older emp-zk. We benchmark activation layers, not end-to-end
networks.

% =====================================================================================
\section{Related work}
Hao et al.~\cite{hao2024scalable} is the framework we build on. Mystique~\cite{weng2021mystique},
Hao et al.'s baseline, provides arithmetic--Boolean conversions and evaluates floating-point
functions with Boolean circuits. zkLLM~\cite{sun2024zkllm} proves arbitrary element-wise functions
with \textsf{tlookup} on GPUs, and its attention protocol zkAttn uses a logarithm. Lu et
al.~\cite{lu2024extensible} reduce non-linear layers to range and exponent relations. zkCNN
covers ReLU and max-pooling~\cite{liu2021zkcnn}. SNARK-based zkML uses lookup
tables~\cite{kang2022scaling}, and Sridhar et al.\ give low-depth provable approximations of
non-polynomial functions~\cite{sridhar2025provable}. In MPC, SIRNN combines digit-wise lookups
for $\exp$ with normalisation and lookups for reciprocal square roots~\cite{rathee2021sirnn}, and
Aly and Smart benchmark scientific functions such as $\sin$ and $\log$ built from range reduction
and polynomials~\cite{aly2019benchmarking}. The VOLE-based ZK backends we use
are surveyed in~\cite{baum2023sok}. We add sound, precise instantiations of $\log$ and periodic
functions \emph{within} a digit-decomposition framework, together with a reusable account of the
pitfalls of the direct instantiations.

\section{Conclusion}
Reusing a digit-wise lookup template for new functions is attractive, but some details are easy to
miss on paper and only show up in an implementation. The template is sound only if every witness is
pinned to the input, and correct only if the recombination is a homomorphism of the function.
Following these two rules gives a logarithm and a periodic-function protocol at 1--1.5 units in the
last place, for 1--2 lookups each, and makes softplus, log-softmax
and snake available in Hao et al.'s framework at \SpeedMin--\SpeedMax$\times$ lower cost than
without lookups. Code, specification, tests and benchmark scripts are released with this paper.

{\small
\bibliographystyle{plainnat}
\bibliography{refs}
}

\end{document}
